\chapter{Constipation Treatment}

\medskip
\noindent\textit{\textbf{Under review.} This chapter is an AI-assisted educational draft; it carries no source citations and is not a substitute for professional medical advice.}

\section*{Definition and Overview}
Treating constipation begins with lifestyle measures and, when needed, laxatives, but the right approach depends on the cause and the person. The goals are to pass soft stools comfortably, to empty the bowel when it is packed (disimpaction), and to establish habits that prevent a recurrence. Most people respond to a combination of more fibre, more fluid, regular toileting, and a short course of laxatives. Persistent or severe constipation may need medical review and specialist care.

\section*{Epidemiology}
Because constipation affects roughly 1 in 7 adults and 1 in 10 children, treatment is among the most common reasons for pharmacy and GP visits. Many people self-treat with laxatives, and prolonged use of stimulant laxatives is common even though it is not usually recommended without advice. Correct treatment that includes dietary change and bowel retraining is more likely to succeed than laxatives alone. In older people, constipation and its complications are a frequent cause of hospitalisation.

\section*{Aetiology and Pathophysiology}
Treatment is directed at the mechanisms causing slow or difficult stool passage:
\begin{itemize}
    \item \textbf{Inadequate fibre and fluid:} dry, hard stool that moves slowly
    \item \textbf{Low activity:} reduced gut movement with immobility
    \item \textbf{Ignoring the urge to go:} the rectum stretches and becomes less sensitive
    \item \textbf{Medicines:} opioids, iron, and other drugs slow the bowel
    \item \textbf{Weak pelvic floor muscles:} difficulty pushing stools out (obstructed defecation)
    \item \textbf{Slow colonic transit:} waste takes longer than usual to reach the rectum
    \item \textbf{Underlying conditions:} hypothyroidism, diabetes, or bowel disease
\end{itemize}

\section*{Clinical Features}
Features that guide treatment include:
\begin{itemize}
    \item Fewer than three bowel movements a week
    \item Hard, dry, or lumpy stools that are painful to pass
    \item Straining and a sense of incomplete emptying
    \item Bloating, wind, and abdominal discomfort
    \item Soiling or leakage in severe or long-standing cases
    \item Bleeding with stools, which can follow hard passage
    \item Need for laxatives already, and how often they are used
\end{itemize}

\section*{Differential Diagnosis}
\begin{itemize}
    \item Irritable bowel syndrome with constipation
    \item Bowel obstruction or a blocked bowel, which needs urgent care
    \item Bowel cancer, especially with new symptoms, bleeding, or weight loss
    \item Hypothyroidism and other metabolic disorders
    \item Medication-induced constipation
    \item Obstructed defecation from pelvic floor problems
    \item Neurological conditions affecting the bowel's nerves
\end{itemize}

\section*{When to Seek Medical Care}
See a doctor if constipation is new and persistent in an older person, does not improve with simple measures, or is accompanied by bleeding, weight loss, or severe pain.

\textbf{Call 000 for:}
\begin{itemize}
    \item Severe abdominal pain with vomiting or a swollen abdomen
    \item Inability to pass wind or stool with worsening pain (possible blockage)
    \item Heavy bleeding from the back passage with faintness
    \item Fever, confusion, or severe dehydration with constipation
\end{itemize}

\section*{Diagnosis and Investigations}
\begin{itemize}
    \item \textbf{History:} stool frequency, consistency, straining, medicines, diet, fluids, and how long symptoms have lasted
    \item \textbf{Examination:} abdominal and rectal examination
    \item \textbf{Blood tests:} thyroid function, blood count, calcium, and glucose if indicated
    \item \textbf{Colonoscopy or flexible sigmoidoscopy:} for older people with new symptoms or warning features such as bleeding
    \item \textbf{Bowel transit studies:} X-rays or markers to measure how long waste takes to travel through the bowel
    \item \textbf{Anorectal physiology and defecography:} to assess pelvic floor function when straining is the main problem
\end{itemize}

\section*{Management}
\begin{itemize}
    \item \textbf{Diet and fluids:} gradually increase fibre to 25--30 grams daily with fruits, vegetables, legumes, and whole grains, and drink plenty of water
    \item \textbf{Exercise:} regular physical activity supports normal bowel function
    \item \textbf{Toilet routine:} sit at the same time each day, ideally after meals, and do not hurry or ignore the urge
    \item \textbf{Bulk-forming laxatives:} psyllium or ispaghula husk are first-line and safe for long-term use
    \item \textbf{Osmotic laxatives:} lactulose or polyethylene glycol soften stools by drawing water into the bowel
    \item \textbf{Stimulant laxatives:} bisacodyl or senna for short-term or occasional use; not recommended daily without medical advice
    \item \textbf{Disimpaction:} for a packed bowel, a doctor may arrange a course of high-dose laxatives
    \item \textbf{Pelvic floor retraining:} biofeedback and physiotherapy help people who strain due to weak or tight pelvic floor muscles
    \item \textbf{Prescribed medicines:} newer prescription agents for severe or chronic constipation unresponsive to usual laxatives
\end{itemize}

\section*{Complications}
\begin{itemize}
    \item Haemorrhoids and anal fissures worsened by straining
    \item Faecal impaction requiring manual or enema treatment
    \item Electrolyte disturbances with prolonged laxative misuse
    \item Dependence on stimulant laxatives
    \item Diarrhoea and cramping from inappropriate laxative doses
    \item Missed serious underlying disease if symptoms are dismissed
\end{itemize}

\section*{Prognosis and Outcomes}
Most cases of constipation improve within days to a few weeks with adequate fluids, fibre, and regular toileting. Chronic constipation often requires ongoing management, but most people find a regimen that keeps them comfortable. Outcomes are best when the cause is identified and treated, when laxatives are used appropriately and not overused, and when pelvic floor or neurological causes receive specialised care. Symptoms that do not respond to standard treatment should be reassessed rather than accepted.

\section*{Prevention and Self-Care}
\begin{itemize}
    \item Eat a varied, fibre-rich diet and increase fibre gradually to avoid bloating
    \item Drink enough fluid, especially as fibre intake rises
    \item Exercise regularly and avoid long periods of sitting
    \item Respond promptly to the urge to use the toilet
    \item Review medicines that cause constipation with a doctor or pharmacist
    \item Avoid regular use of stimulant laxatives without advice
    \item Seek medical review if symptoms persist despite these measures
\end{itemize}
